\pdfoutput=1
\documentclass[11pt]{article}
\usepackage{emnlp2023}
\usepackage{times}
\usepackage{latexsym}
\usepackage[T1]{fontenc}
\usepackage[utf8]{inputenc}
\usepackage{microtype}
\usepackage{booktabs}
\usepackage{graphicx}
\usepackage{amsmath}

\title{Acknowledged Is Not Done: Reward Observability, Evidence Use and\\Recovery Training for Tool Agents with Silent Write Failures}

\author{Atharv Sharma \\
  BITS Pilani, Goa Campus \\
  \texttt{atharvgaur591@gmail.com}}

\newcommand{\RV}{R$_V$}
\newcommand{\RA}{R$_A$}

\begin{document}
\maketitle

\begin{abstract}
When a tool agent writes to an external system, the tool's acknowledgement can be identical whether the write committed, half-committed, or did nothing. We study GRPO-trained tool agents in a procurement simulator built around this failure, with an evaluator that sees the authoritative database and a policy that sees only tool replies. Pre-registered experiments on a 3.35B Aya model show: (1) a state reward and an acknowledgement reward of identical binary form produce agents that do and do not read back after phantom writes in every one of five seeds (1.00 vs 0.00), with the acknowledgement-trained agent behaving exactly as its untrained warm start; (2) once the authoritative record is in context, every model acts correctly on whether a payment exists, but only state-trained agents refund a partial payment, and only when payment is the first step of the workflow; (3) every model linearly encodes partial versus complete payment, and training changes how that evidence moves the next action rather than whether it is represented. We then test a recovery-start curriculum, a known technique: it raises verified completion by +0.21 on fresh tasks in five of five seeds, but a pre-registered class-matched control attributes about half of the gain to fault coverage, and the remaining start-state effect (+0.09, three seeds) comes with slightly more irreversible violations. Gains do not transfer to two sibling domains. Every number traces to a result file through a public claims ledger.
\end{abstract}

\section{Introduction}

A payment API returns \texttt{\{"status": "ok"\}}. The payment may have been stored, stored in part, stored twice, or not at all. An agent that reports success because the tool said so is wrong in exactly the cases that matter, and retrying blindly can pay twice. We call such failures \emph{silent}: the agent's observation does not distinguish them; only the authoritative state does.

We ask three questions in a controlled simulator where silent failures are the rule rather than the exception. \textbf{Acquisition}: does the information a reward observes decide whether a trained agent checks the authoritative state after a write? \textbf{Use}: once the evidence is in context, does the agent act on it correctly? \textbf{Training}: does a recovery-start curriculum improve recovery, and how much of its gain comes from starting at recovery states rather than from seeing more fault types?

\paragraph{Contributions.} (i) ProcureSim, an environment in which phantom, partial, lost-acknowledgement, duplicate and never-committing writes return acknowledgements byte-identical to real outcomes, with a state-only verifier checked against an independent reference implementation and hand-written golden traces (\S\ref{sec:setup}). (ii) Pre-registered measurements of acquisition, evidence use and representation under two rewards that differ only in what they observe (\S\ref{sec:acq}--\S\ref{sec:mech}). (iii) A decomposition of a recovery-start curriculum's gain into fault coverage and start state, with its safety trade-off (\S\ref{sec:rsc}). (iv) Surface, language, workflow and fault-combination shifts of the benchmark, and a test of a forked counterfactual credit estimator in this setting (\S\ref{sec:shift}). \textbf{Non-contributions}: we do not propose a new optimiser (the curriculum and the credit estimator adapt published methods; \S\ref{sec:related}), and we claim nothing about real deployments.

\section{Related work}\label{sec:related}

\paragraph{Tool reliability and false success.} Agents claim completion that the environment contradicts \citep{smyth2026overclaim}; Proof-of-Use studies tool calls that do not ground answers \citep{ma2026proofofuse}. Bench2Robust trains recovery from explicit and silently corrupted tool outputs but assumes a retry cannot duplicate a write \citep{chen2026bench2robust}. Our failures are silent at the write itself: the acknowledgement is identical, and retries can duplicate.

\paragraph{Recovery and restart training.} Restarting GRPO from failure contexts is established: Fission-GRPO turns failed calls into corrective prompts \citep{zhang2026fission}, reverse curricula start RL from late points of a trajectory \citep{xi2024r3}, and ReflectPO and Depth Before Breadth restart from recovery anchors and intermediate contexts \citep{li2026reflectpo,du2026depth}. Our recovery-start curriculum belongs to this family; our contribution is its decomposition in a silent-failure setting.

\paragraph{Credit at forked states.} Forking environment snapshots to estimate step credit appears in CRR \citep{li2026crr} and SURE-Act \citep{liu2026sureact}; action-conditioned baselines in BiPACE \citep{wang2026bipace}, GiGPO \citep{feng2025gigpo} and COMA \citep{foerster2017coma}. Constrained GRPO and its calibration \citep{girgis2026cgrpo,han2026budget} handle separate cost terms. Curricula that oversample hidden states can bias the learned policy \citep{jiang2022samplr}.

\section{Setting}\label{sec:setup}

\paragraph{ProcureSim.} An agent processes an approved purchase order in one of three workflows (order then pay then post; receive then pay then post; pay then post) using five read tools, six write tools and a final report. Faults are applied at one write chokepoint and are invisible in the acknowledgement: \emph{phantom} (ok, nothing stored), \emph{partial} (ok, a fraction $q$ stored), \emph{lost acknowledgement} (stored, timeout returned), \emph{timeout before commit}, \emph{duplicate delivery} (stored twice, ok returned) and \emph{never commits} (ok on every attempt, nothing stored). A same-key retry after a lost acknowledgement is answered from an idempotency store; a partial payment can only be repaired by refunding it and paying again under a new key. Table~\ref{tab:truth} gives the reward each behaviour receives.

\paragraph{Rewards and metric.} \RV{}$=1$ iff every requested step verifies in the final database. \RA{}$=1$ iff the last acknowledgement of every requested write was ok. Both are binary and ignore the report. The evaluation metric is verified safe completion (VSC): every step verified, no collateral or unrequested row, and a \emph{completed} report. False completion is a \emph{completed} report when the goal is not met; a violation is a duplicate payment or receipt, a collateral row, or a wrong amount or payee.

\paragraph{Training.} Tiny Aya Global (3.35B) with LoRA (rank 16) is warm-started by supervised fine-tuning on scripted episodes (a verifying oracle in 30\% of them, an acknowledgement-trusting policy otherwise), then trained for 200 steps with TRL's GRPOTrainer: 4 prompts $\times$ 8 rollouts per step, learning rate $2\times10^{-5}$, KL coefficient 0.02, group-normalised sequence advantages, token-level DAPO loss aggregation and sequence-level importance correction (the trainer defaults; Appendix~\ref{app:train}).

\begin{table}[t]
\centering
\resizebox{\columnwidth}{!}{% generated by scripts/paper_tables.py; do not edit
\begin{tabular}{lcccc}
\toprule
Fault & proceed & blind retry & read, retry & read, stop \\
\midrule
clean & 1/1 & 1/1 & 1/1 & 1/1 \\
phantom & 0/1 & 0/1 & 1/1 & 0/1 \\
partial & 0/1 & 0/1 & 0/1 & 0/1 \\
lost ack & 1/0 & 1/1 & 1/0 & 1/0 \\
timeout before commit & 0/0 & 1/1 & 1/1 & 0/0 \\
never commits & 0/1 & 0/1 & 0/1 & 0/1 \\
\bottomrule
\end{tabular}
}
\caption{\RV{}/\RA{} for scripted behaviours under each fault (registered reward truth table). The rewards disagree exactly where reading changes the outcome.}
\label{tab:truth}
\end{table}

\section{Acquisition: what the reward observes}\label{sec:acq}

We compared \RV{} and \RA{} in a pre-registered design with five paired seeds, each evaluated once on 300 untouched tasks (Table~\ref{tab:a12}). After a phantom write, the \RV{} agent reads the payment record in every case and every seed (1.00); the \RA{} agent never does (0.00). VSC rises by +0.27 to +0.32 per seed (exact sign-flip $p=1/32$ one-sided; Holm-adjusted 0.0625 across the two primary metrics). The \RA{} agent's greedy tool calls equal the untrained warm start's on 500 of 500 confirmatory tasks in each seed, so this contrast is between a state-trained agent and one that training did not move. Our registered mechanism prediction (per-class sign predictions) failed and is reported as such. A pre-registered control that removed from \RA{}'s training pool the classes where reading lowers \RA{} left greedy read-back at 0/153 in all six runs, while the sampled read rate stopped declining: pool composition explains the drift away from reading, not its absence.

\begin{table*}[t]
\centering
\resizebox{0.86\textwidth}{!}{% generated by scripts/paper_tables.py; do not edit
\begin{tabular}{lcccccc}
\toprule
 & \multicolumn{2}{c}{Phantom read-back} & \multicolumn{2}{c}{VSC} & \multicolumn{2}{c}{False success} \\
\cmidrule(lr){2-3}\cmidrule(lr){4-5}\cmidrule(lr){6-7}
Seed & R$_V$ & R$_A$ & R$_V$ & R$_A$ & R$_V$ & R$_A$ \\
\midrule
warm start & \multicolumn{2}{c}{0.00} & \multicolumn{2}{c}{0.47} & \multicolumn{2}{c}{0.53} \\
\midrule
0$^\dagger$ & 1.00 (153/153) & 0.00 (0/153) & 0.783 & 0.470 & 0.060 & 0.530 \\
1 & 1.00 (153/153) & 0.00 (0/153) & 0.793 & 0.470 & 0.050 & 0.530 \\
2$^\ddagger$ & 1.00 (153/153) & 0.00 (0/153) & 0.783 & 0.470 & 0.060 & 0.530 \\
3 & 1.00 (153/153) & 0.00 (0/153) & 0.747 & 0.470 & 0.097 & 0.530 \\
4 & 1.00 (153/153) & 0.00 (0/153) & 0.743 & 0.470 & 0.100 & 0.530 \\
\midrule
\multicolumn{7}{l}{VSC difference per seed: +0.313, +0.323, +0.313, +0.277, +0.273; exact sign-flip $p$ = 1/32 one-sided, 2/32 two-sided} \\
\bottomrule
\end{tabular}
}
\caption{A12 confirmatory results per seed. $^\dagger$Seed 0 ran under the original protocol (stop rule fired; deviation D-A12-1). $^\ddagger$Resumed from step 175 after a crash (D-A12-2).}
\label{tab:a12}
\end{table*}

\section{Use: acting on evidence}\label{sec:use}

We force the authoritative read into the agent's context after a pay acknowledgement, under five real execution states (pre-registered). Every policy, trained or not, re-pays when the record is empty and proceeds when it is full (branch-balanced accuracy 1.00). Only \RV{} agents refund a partial payment, and a second registered experiment shows they do so at unseen ratios (1/4, 3/4) as often as at the 1/2 seen in training: 0.25--0.58 in four seeds and 0 in the fifth, against 0.00 for \RA{} and 0.035 for the warm start. Split by workflow, the picture changes: in seeds 0--2 the \RV{} agent repairs 99--100\% of partial payments when payment is the workflow's first step and 0\% when it follows a goods receipt. The amount comparison is real but bound to the context in which it was learned. No model refunds a duplicate delivery.

\section{Representation}\label{sec:mech}

At the decision point after the forced read, a linear probe separates partial from complete payment with AUROC 1.00 at every checkpoint of every model, including the warm start, whether refitted or frozen from the warm start. What differs is the action margin $M=\log P(\text{refund or escalate})-\log P(\text{post})$: its partial-minus-complete gap grows under \RV{} (step 200: +2.9, +3.9, +5.1 nats) and stays near the warm start under \RA{} (+2.0 vs +2.4). Projecting out the dev-fitted partial-minus-complete direction at the registered layer lowers \RV{}'s partial-branch margin more than any random direction in five of five seeds, meeting our causal criterion; the same ablation also lowers the margin on the empty-record branch and meets the criterion for two \RA{} seeds, so it is not specific to the \RV{} behaviour. Single-position activation patching did not exceed random controls in any seed. Training changes how represented evidence drives the action, not whether it is represented; we do not claim a circuit.

\section{A recovery-start curriculum, decomposed}\label{sec:rsc}

\RV{}'s remaining failures are partial payments, duplicate deliveries and false reports, and GRPO rarely samples the repair from a full episode. A recovery-start curriculum (RSC) replays a scripted prefix ending in a faulted state, so that each group of eight rollouts starts there; the update rule and the reward (\RV{}) are unchanged. The pool is half \RV{}'s original tasks and half recovery starts, which also add fault types \RV{} never trained on (duplicate delivery, partial ratios 1/4 and 3/4).

\paragraph{Against \RV{} (recipe comparison).} On fresh pre-registered splits, RSC raises VSC by +0.21 on recovery starts and on full episodes, in five of five seeds (Table~\ref{tab:s3}), lowers false completion by 0.21 and violations by 0.12, and leaves honest failure unchanged (no collapse to refusal).

\paragraph{Against a class-matched control (start-state comparison).} Registered before any S3 output was read, the control trains \RV{} on full episodes carrying exactly the recovery starts' faults, prompts and worlds. Against it, starting at recovery states adds +0.09 VSC on recovery starts and +0.22 on full-episode partial and duplicate tasks (positive in each of three seeds; sign-flip resolution 1/8; intervals are task bootstraps conditional on the trained models). It also adds violations (+0.03 on that subset) and is worse on duplicate stress ($-0.04$): the stress gain over \RV{} comes from fault coverage. The start-state gain is concentrated in the order-then-pay workflow (Table~\ref{tab:wf}).

\begin{table}[t]
\centering
\resizebox{\columnwidth}{!}{% generated by scripts/paper_tables.py; do not edit
\begin{tabular}{llccc}
\toprule
Contrast & Split & $\Delta$VSC & $\Delta$violation & Seeds $>0$ (VSC) \\
\midrule
RSC $-$ R$_V$ (recipe) & recovery starts & +0.21 [+0.17, +0.25] & -0.11 [-0.15, -0.08] & 5/5 \\
RSC $-$ R$_V$ (recipe) & full episodes & +0.21 [+0.17, +0.25] & -0.12 [-0.15, -0.08] & 5/5 \\
RSC $-$ R$_V$ (recipe) & stress & +0.39 [+0.30, +0.48] & -0.39 [-0.48, -0.30] & 5/5 \\
RSC $-$ matched (start state) & recovery starts & +0.09 [+0.06, +0.11] & +0.01 [-0.00, +0.02] & 3/3 \\
RSC $-$ matched (start state) & partial + duplicate & +0.22 [+0.15, +0.29] & +0.03 [+0.00, +0.05] & 3/3 \\
RSC $-$ matched (start state) & stress & -0.04 [-0.07, -0.01] & +0.04 [+0.01, +0.07] & 1/3 \\
RSC $-$ matched, 5 seeds (Holm $p$=0.036) & recovery starts & +0.11 [+0.08, +0.13] & -- & 5/5 \\
RSC $-$ matched, 5 seeds (Holm $p$=0.111) & partial (full episodes) & +0.38 [+0.32, +0.45] & -- & 4/5 \\
RSC $-$ matched, 5 seeds (Holm $p$=0.337) & duplicate (full episodes) & +0.13 [+0.09, +0.17] & -- & 2/5 \\
\bottomrule
\end{tabular}
}
\caption{Registered S3 contrasts: seed-averaged paired difference [95\% CI]. Five seeds against \RV{}, three against the class-matched control.}
\label{tab:s3}
\end{table}

\begin{table}[t]
\centering
\resizebox{\columnwidth}{!}{% generated by scripts/paper_tables.py; do not edit
\begin{tabular}{lccc}
\toprule
Model & order$\to$pay & receive$\to$pay & pay first \\
\midrule
GRPO R$_V$ & 17/111 (0.15) & 16/102 (0.16) & 63/117 (0.54) \\
class-matched & 62/111 (0.56) & 63/102 (0.62) & 54/117 (0.46) \\
RSC & 109/111 (0.98) & 61/102 (0.60) & 82/117 (0.70) \\
\bottomrule
\end{tabular}
}
\caption{VSC on full-episode partial and duplicate tasks by workflow (seeds 0--2 pooled).}
\label{tab:wf}
\end{table}

\section{Shifts and a counterfactual credit test}\label{sec:shift}

% PENDING: S4-EVAL (PREREGISTRATION_S4_EVAL.md) and FCA (docs/research/algorithm/FCA_GRPO.md). Fill only from results/s4/.
\textit{[Pending: completed only from the registered result files.]}

In two separately implemented sibling domains (IT operations, healthcare administration), neither \RV{} nor RSC improves verified completion over the warm start, and both trained models make more collateral writes than the warm start in healthcare administration. The \RV{}-over-\RA{} read-back gap transfers to healthcare administration (+0.37 [0.18, 0.58]) but not to IT operations, where the warm start already reads back.

\section{Discussion}

The information a reward observes decides whether an agent checks; this is the clearest result, and it replicates across five seeds. Checking does not imply using: reading is followed by correct action on present-or-absent evidence in every model, but quantitative repair is learned narrowly, by workflow. Starting training at recovery states helps beyond seeing the faults, but less than the recipe comparison suggests, and not for free. None of this transfers to other domains without training there.

\section*{Limitations}

One model family (Tiny Aya Global, LoRA rank 16) and one authored simulator; the sibling domains share code and authorship. Five seeds for primary comparisons and three for the start-state control; task-level intervals are conditional on the trained models. The recovery-start pool differs from \RV{}'s in fault coverage, which the class-matched control isolates only for three seeds. The \RA{} arm saturates and stays at warm-start behaviour, so A12 compares a trained agent with an untrained one. Probes and ablations do not localise a mechanism. All results are in simulation.

\section*{Ethics Statement}

The simulator processes synthetic purchase orders; no real payment, person or organisation is involved. Agents that act on external systems can cause irreversible harm; our results argue for checking authoritative state rather than trusting tool acknowledgements.

\bibliography{refs}
\bibliographystyle{acl_natbib}

\appendix
\section{Environment and tools}
\label{app:env}

ProcureSim (\texttt{actionproof/integrity/sim/}) exposes twelve tools. Reads: \texttt{get\_purchase\_order}, \texttt{find\_supplier\_order}, \texttt{list\_payments}, \texttt{list\_goods\_receipts}, \texttt{get\_journal\_entries}. Writes: \texttt{place\_supplier\_order}, \texttt{pay\_supplier}, \texttt{record\_goods\_receipt}, \texttt{post\_journal\_entry}, \texttt{refund\_payment}, \texttt{escalate\_to\_human}; and \texttt{submit\_report} (status and summary). Every write passes one chokepoint that applies the task's fault plan to its attempt index and records whether the write committed nothing, part, all, or twice. The verifier (\texttt{truth.py}) reads only the final state; an independently written reference verifier and 91 independently written golden traces agree with it (\texttt{tests/golden/}). The specification is \texttt{docs/stage2/VERIFIER\_SPEC.md}.

\section{Splits}
\label{app:splits}

All splits are generated from fixed seeds by \texttt{actionproof/bench/generator\_proc.py}; task and purchase-order ids are disjoint across splits (tested), and the prompt never depends on the fault plan (tested). Training: \texttt{proc\_train\_b} (400 tasks, rebalanced towards classes where the rewards disagree; registered addendum 1). Development: \texttt{proc\_dev} (200), \texttt{proc\_val} (100). Confirmatory, each frozen model run once: \texttt{proc\_conf} (300; clean 20\%, phantom 25\%, partial 10\%, lost acknowledgement 15\%, timeout before commit 15\%, never commits 15\%), \texttt{proc\_combo} (100 unseen fault pairs), \texttt{proc\_stress} (100: duplicate delivery; lost acknowledgement followed by a lagging read). A14: \texttt{proc\_cf} (250, generated after its registration). Hashes: \texttt{results/a12*\_hashes.json}, \texttt{results/a14b\_hashes.json}; card: \texttt{docs/stage2/BENCHMARK\_CARD.md}.

\section{Training configuration}
\label{app:train}

Warm start: supervised fine-tuning (2 epochs, lr $10^{-4}$, LoRA rank 16) on 240 scripted episodes, each from the verifying oracle with probability 0.30 or the acknowledgement-trusting policy otherwise. GRPO (TRL with AgentTune's batched rollouts): 200 steps, 4 tasks $\times$ 8 rollouts per step, lr $2\times10^{-5}$, KL coefficient 0.02, temperature 1, at most 16 turns and 256 generated tokens per turn, LoRA rank 16; one H200 per run (about 1 h 50 min including development evaluations). Seeds 0--4 per reward; seed 0 under the original registration, seeds 1--4 under addendum 2, registered before they started.

\section{Statistics}
\label{app:stats}

Unit: task. Per seed: paired task bootstrap (10{,}000 resamples) and exact McNemar test. Across seeds: per-seed estimates, direction consistency and the exact sign-flip test (smallest attainable $p$: 1/32 one-sided and 2/32 two-sided at five seeds); a two-level bootstrap that resamples seeds and shares task draws across seeds is reported as descriptive. Primary family (phantom read-back, VSC): Holm. A14: task bootstrap; categories with a smallest effect of 0.10 (evidence-conditioned, insensitive, inverted, inconclusive); contrasts are \emph{different} if the 95\% interval excludes 0 and \emph{equivalent} if the 90\% interval lies within $\pm$0.10. Code: \texttt{scripts/s2\_stats.py}, \texttt{scripts/a14\_stats.py}. S3 (\texttt{scripts/s3\_stats.py}, \texttt{s3\_stats\_addendum.py}): per-seed paired task bootstrap and a seed-averaged bootstrap with one task draw shared by all seeds; exact sign-flip across seeds. Pairing RSC seed $s$ with \RV{} seed $s$ shares no randomness (different pools and code paths), so we also report an unpaired seed permutation over all assignments of per-seed means (smallest $p$: 1/252 for five runs per arm, 1/20 for three); S4-EVAL uses it as the registered decision statistic (\texttt{scripts/s4\_stats.py}).

\section{Track 3 metrics}
\label{app:track3}

\begin{table}[h]
\centering
\resizebox{\columnwidth}{!}{% generated by scripts/paper_tables.py; do not edit
\begin{tabular}{lccccc}
\toprule
Model & Verified & Claimed & Tool-call validity & Calls/ep. & Latency s (p50/p95) \\
\midrule
warm start & 0.47 & 1.00 & 1.00 & 3.9 & 2.41/3.05 \\
R$_A$ seed 0 & 0.47 & 1.00 & 1.00 & 3.9 & 2.40/3.05 \\
R$_V$ seed 0 & 0.78 & 0.84 & 1.00 & 7.3 & 3.62/4.68 \\
\bottomrule
\end{tabular}
}
\caption{\texttt{proc\_conf}: verified completion next to claimed success (share of episodes reported \texttt{completed}); latency is batch-1 seconds per episode on 40 \texttt{proc\_dev} tasks (L40S). The gateway adds about 0.2\,ms per tool call (in process).}
\end{table}

\section{Gateway factorial (A11, e-commerce)}
\label{app:a11}

\begin{table}[h]
\centering
\resizebox{\columnwidth}{!}{% generated by scripts/paper_tables.py; do not edit
\begin{tabular}{lcccccc}
\toprule
Policy & Passthrough & Enforce & Enf.+Inform & $\Delta_\text{runtime}$ & $\Delta_\text{information}$ & $I$ \\
\midrule
warm start & 0.43 & 0.62 & 0.87 & +0.19 & +0.25 & -- \\
PCSTR-d s0 & 0.65 & 0.94 & 0.94 & +0.29 & +0.00 & +0.107 \\
PCSTR-d s1 & 0.73 & 0.94 & 0.94 & +0.21 & +0.00 & +0.027 \\
ack s0 & 0.66 & 0.81 & 0.87 & +0.15 & +0.06 & -0.040 \\
ack s1 & 0.62 & 0.80 & 0.87 & +0.18 & +0.07 & -0.007 \\
tool trust s0 & 0.64 & 0.82 & 0.87 & +0.18 & +0.06 & -0.007 \\
\bottomrule
\end{tabular}
}
\caption{A11 (registered; Stage 1 e-commerce environment, 300 tasks): system-level verified completion per gateway mode. $I$: registered interaction (policy effect under Enforce minus under Passthrough); it excludes zero for PCSTR-d seed 0 only, so the registered claim, which required both seeds, is not made. The runtime check ($\Delta_\text{runtime}>0$ for every policy) passes.}
\end{table}

\section{Contract coverage (A13) and defences (H7)}
\label{app:a13}

The registered A13 test (one-sided Fisher, adversarial vs represented violations) treated repeated deterministic rows as independent samples; we do not use its $p$-value (deviation D-A13-1). As finite-set results: the contract accepts none of the 194 single-fault and 479 two-fault violating configurations it was written for (an exhaustive enumeration; no configuration varies across the 200 worlds), and 14 of 161 independently written violations (8 of 151 unambiguous), mostly refunds and master data; every false accept is listed in \texttt{results/a13/h5\_finite.json}. On identical fault traces, idempotency keys alone recover no phantom write and verify-before-retry recovers all; this table is deterministic and true by construction (\texttt{results/a11/h7\_matrix.md}).

\section{Stage 1 results}
\label{app:stage1}

Stage 1 (Nov 1--15, 2026, e-commerce environment) is prior work. Its registered confirmatory result: PCSTR-d (a state-grounded reward with report terms) reached VSC 0.8425 against 0.8275 for the best baseline, +1.5 points [0.6, 2.6] on 400 tasks; both PCSTR-d seeds give identical outcomes. Not registered, and between rewards of different form: a tool-trusting reward reached 0.640. With rewards of identical form on a development split, sampled refund read-back was 0.25 for the warm start, 0.71--0.72 under state rewards and 0.10 under the acknowledgement reward. An LLM judge (Aya Expanse 32B, AuditKIT prompt) given the receipt ledger still accepted 71\% of false completions.

\section{Human audit (H4)}
\label{app:h4}

Registered as a between-subjects audit of 60 episodes by one annotator (\texttt{PREREGISTRATION\_A11.md} \S9). On failures that the transcript does not show, the annotator accepted 8 of 10 as completed from the transcript alone and 1 of 10 with the receipt ledger; the registered prediction ($\geq$0.5 and $\leq$0.2) is met and the falsification threshold ($<$0.3) is not reached. Descriptively, failures the transcript does show were accepted 10 of 10 from the transcript and 0 of 10 with the ledger. One annotator and ten items per cell: wide intervals (\texttt{results/h4/h4\_summary.json}).


\section{Recovery-start curriculum and class-matched control}
\label{app:rsc}

\textbf{Recovery starts} (\texttt{actionproof/stage3/tasks.py}). A task carries a scripted prefix built with the oracle's own argument and key conventions, ending (i) at the acknowledgement of a faulted write (phantom; partial at ratio 1/4, 1/2 or 3/4; lost acknowledgement; duplicate delivery on pay; complete; never commits), (ii) after its authoritative read, or (iii) after the whole workflow is complete. The rollout (\texttt{actionproof/stage3/rollout.py}) replays the prefix into the simulator and the conversation; the replayed turns are prompt (no loss). Training pool: 200 \texttt{proc\_train\_b} tasks plus 200 recovery starts; trainer, reward (\RV{}), steps and rollouts are those of A12.
\textbf{Development selection} (\texttt{docs/stage3/RESEARCH\_REPORT.md} \S6): four 100-step pilots (a lexicographic safety-first reward, recovery starts with \RV{}, both, and a Lagrangian penalty) and two 200-step mixes (50\% and 25\% recovery starts, two seeds) were compared on development splits only; the fixed rule selected 50\% recovery starts with \RV{}. The safety-first and Lagrangian rewards gave no development gain and were not carried forward.
\textbf{Correction.} The S3 registration stated that only the start states changed; addendum 1, registered before any S3 output was read, corrects this (the pool also adds fault types) and adds the class-matched control: the same 200 \texttt{proc\_train\_b} tasks plus 200 full episodes built from the recovery starts themselves (same world, prompt and fault; prefix removed).

\section{Benchmark shifts (S4-EVAL)}
\label{app:shift}

Registered in \texttt{PREREGISTRATION\_S4\_EVAL.md}, generated by \texttt{scripts/s4\_make\_splits.py}: the same 150 tasks with the request in English, Hindi, Hinglish or reworded English (paired); 150 tasks behind a different API surface (all twelve tools and their arguments renamed, acknowledgements \texttt{\{"result": "ACCEPTED"\}}, timeouts \texttt{deadline\_exceeded}; canonical names rejected; same actions give the same final state and score, tested); 100 tasks of an unseen four-step workflow (order, receive, pay, post); 100 tasks with two faults on two steps. Scripted oracle VSC 0.81--1.00 (below 1 exactly on never-commit tasks).

\textbf{Results} (Table~\ref{tab:s4}; \texttt{results/s4/s4\_summary.json}). The registered rule, amended before any output existed (addendum 1: Holm over the five-seed sign-flip could not reach 0.05, so the decision statistic became the unpaired seed permutation over all 252 assignments), supports G2 (Hindi and Hinglish) and G4 (two faults) and rejects G1 (renamed API) and G3 (unseen workflow). Against the class-matched control (three seeds, descriptive), RSC is lower on the renamed API and the unseen workflow and within a few points elsewhere. Diagnostics (\texttt{scripts/s4\_diagnose.py}) are in the main text.

\begin{table}[h]
\centering
\resizebox{\columnwidth}{!}{% generated by scripts/paper_tables.py; do not edit
\begin{tabular}{lcccccc}
\toprule
Split & warm & R$_V$ & RSC & matched & RSC $-$ R$_V$ [95\% CI] & Seeds $>0$ \\
\midrule
English & 0.33 & 0.70 & 0.83 & 0.78 & +0.12 [+0.08, +0.17] & 5/5 \\
Hindi (G2) & 0.35 & 0.68 & 0.81 & 0.78 & +0.13 [+0.08, +0.18] & 5/5 \\
Hinglish (G2) & 0.33 & 0.66 & 0.74 & 0.74 & +0.07 [+0.04, +0.11] & 5/5 \\
reworded English & 0.53 & 0.71 & 0.83 & 0.78 & +0.12 [+0.08, +0.17] & 5/5 \\
renamed API (G1) & 0.07 & 0.21 & 0.20 & 0.29 & -0.02 [-0.04, +0.01] & 3/5 \\
unseen workflow (G3) & 0.12 & 0.03 & 0.05 & 0.11 & +0.02 [+0.00, +0.03] & 2/5 \\
two faults (G4) & 0.21 & 0.59 & 0.66 & 0.63 & +0.07 [+0.03, +0.12] & 5/5 \\
\bottomrule
\end{tabular}
}
\caption{S4-EVAL: mean VSC per arm (5, 5 and 3 runs) and the registered RSC $-$ \RV{} contrast (seed-averaged paired task bootstrap).}
\label{tab:s4}
\end{table}

\section{Forked counterfactual credit (FCA)}
\label{app:fca}

Specification: \texttt{docs/research/algorithm/FCA\_GRPO.md}. At a recovery-start decision state, four of a group's eight slots are forced branches, one per action class (proceed, retry, compensate, escalate, or read), whose first turn is a canonical action computed from the visible conversation only; the rest of each branch is sampled. Per-token advantages: on-policy tokens $(r_i-\hat V)/\sigma$; a forced branch's action tokens $(\hat Q(a)-\hat V)/\sigma$, continuation tokens 0, where $\hat V$ is the on-policy mean, $\hat Q(a)$ the mean reward of completions whose first decisive action is class $a$, and $\sigma$ the group's reward standard deviation. Forced tokens are logged with sampling log-probability 0, so TRL's sequence-level importance correction weights a forced branch by $\pi_{\text{old}}(a)$, giving the all-action gradient term $\nabla\pi(a)(\hat Q(a)-\hat V)$. This adapts published estimators (COMA, BiPACE, GiGPO, CRR) to simulator forks; it is not a new optimiser. Kill criteria were fixed before training. \textbf{Result} (\texttt{results/s4/fca\_dev\_summary.json}): development VSC on recovery starts 0.640 and 0.780 against 0.807 and 0.807 for recovery-start GRPO (seeds 0, 1), failing the first criterion; KL stayed within 1.08 times the baseline's. The losses are on partial payments (seed 0) and duplicates. FCA was dropped without a test evaluation (\texttt{docs/research/algorithm/FCA\_RESULT.md}).

\section{Corrections}
\label{app:corr}

Recorded in \texttt{docs/CONSISTENCY\_AUDIT.md} (18 items) and \texttt{docs/stage2/DECISIONS.md}; registrations are never edited. The most consequential: the S3 recipe-versus-start-state correction (above); the safety trade-off against the class-matched control; the stress gain attributed to coverage; the training rollout environment seed is constant (0) in all runs because of how it was derived, which does not affect validity; the trainer's loss is DAPO-style token aggregation.

\section{Deviations}
\label{app:dev}

D-A12-1: the registered stop rule (zero-variance groups $>$0.80 at steps 40--50) fired for \RA{} seed 0 (0.864, reward 0.97); the run was continued and is labelled. Addendum 2 (registered before seeds 1--4) keeps the rule as a diagnostic for a saturated arm. D-A12-2: \RV{} seed 2 crashed at step 199 (tokenizer error in the rollout loop) and was resumed from its step-175 checkpoint. D-A13-1: see Appendix~\ref{app:a13}. D-K1-1: first K1 launches failed before training (port clash; missing input); relaunched. D-A17-1: the seed-averaged A17 interval was computed in a separate script because the frozen one lacked it. D-A16-1: CircuitKIT did not run (dependency error) and was dropped. D-S4E-1: the S4-EVAL addendum states its registration time as about 20:00; its commit time, 19:38, is the registration time. D-S4E-2: the frozen S4 statistics script paired language splits by task id, which carry the split name, and stopped before writing output; a wrapper renames the ids to their shared index and runs it unchanged. D-ID1-1: the ID1 tasks lacked the flag that the evaluation script requires for a confirmatory split; identical copies with the flag were run. Every decision, with its time, is in \texttt{docs/stage2/DECISIONS.md}.

\section{Compute, reproduction and declarations}
\label{app:repro}

Compute: Hugging Face Jobs (H200 for training, L40S for evaluation) for A11--A17 and K1, with every job and its cost in \texttt{docs/stage2/HF\_CREDITS.md}; an institutional Slurm cluster (two H200 GPUs; one A100) for S3, S4, ID1 and S3 addendum 2, with task scripts in \texttt{slurm/tasks/}. Code, pre-registrations, results and the claim ledger that resolves every number in this paper are in the repository (\url{https://github.com/atharvgaur1845/actionproof-v2}); \texttt{scripts/s2\_numcheck.py} re-reads each number from its file.

\textbf{Prior work.} ActionProof Stage 1 (Nov 1--15) built the e-commerce environment, its rewards and runs, and the judge audit. Stage 2 (Nov 20--25, including the internal research phases labelled Stage 3 and Stage 4 in the repository) built ProcureSim, the two rewards studied here, the gateway, the A11--A17, K1, S3 and S4 registrations, all training and evaluation reported here, the product application, and this paper. The models are Cohere's; AgentTune and AuditKIT are the hackathon's toolkits.

\textbf{LLM use.} Claude (Anthropic) was used as a coding and writing assistant. Every number in this paper was produced by code in the repository and is checked against its result file.


\end{document}
